\rtmtight
\begin{tblr}{width=\textwidth,colspec={|Q[c,m,2.1cm]|X[l,m]|},hlines,vlines,hspan=minimal,rowsep=1.5pt,colsep=3pt,row{1}={bg=headgray,font=\bfseries}}
\SetCell{c,m}\hfil\logo\hfil & \textbf{POLITEKNIK NEGERI MALANG}\par\textbf{JURUSAN TEKNOLOGI INFORMASI}\par\textbf{PROGRAM STUDI: D4 TEKNIK INFORMATIKA} \\
\SetCell[c=2]{l} \textbf{RENCANA TUGAS MAHASISWA (RTM) DAN RUBRIK PENILAIAN} & \\
Nama Mata Kuliah & Workshop Teknologi Terapan \\
Kode Mata Kuliah & RTI256205 \quad \textbf{sks / Jam perminggu:} 6 SKS/12 jam \quad \textbf{Semester:} 6 \\
Dosen Pengampu & Tim Pengajar Program Studi D4 TI \\
\end{tblr}
\assessmentblock{Rencana Workshop Industri}{Case Method (CM) -- perencanaan terstruktur}{SCPMK0803-05403, SCPMK0802-05402}{Mahasiswa menyusun rencana workshop industri yang komprehensif mencakup WBS, timeline sprint, target deliverable, dan rencana pengembangan produk berbasis kebutuhan industri.}{Menganalisis kebutuhan industri, menyusun WBS dan timeline, merancang backlog awal, mendistribusikan peran tim, dan mendokumentasikan rencana.}{Dokumen rencana workshop (WBS, timeline, backlog), presentasi rencana kepada stakeholder, dan/atau laporan perencanaan.}{Kelengkapan WBS dan timeline workshop & 6 & WBS mencakup semua deliverable sprint, timeline realistis dan terukur, serta milestone dikonfirmasi stakeholder industri. & WBS dan timeline tersedia tetapi beberapa deliverable atau milestone belum terperinci atau dikonfirmasi. & WBS tidak lengkap, timeline tidak realistis, atau rencana tidak mencerminkan kebutuhan industri. \\ \\
Kualitas analisis kebutuhan dan target deliverable & 5 & Kebutuhan industri dianalisis secara mendalam, target deliverable terdefinisi jelas, dan backlog awal mencerminkan prioritas bisnis. & Analisis kebutuhan cukup, tetapi target deliverable atau prioritas backlog masih perlu penyempurnaan. & Analisis kebutuhan dangkal, target tidak terukur, atau backlog tidak mencerminkan prioritas industri. \\ \\
Distribusi peran tim dan stakeholder plan & 4 & Peran tim terdistribusi sesuai kompetensi, jalur komunikasi stakeholder jelas, dan rencana eskalasi terdokumentasi. & Distribusi peran dan rencana komunikasi ada tetapi beberapa aspek koordinasi belum jelas. & Distribusi peran tidak jelas, tidak ada rencana komunikasi stakeholder, atau perencanaan tim tidak operasional. \\}

\assessmentblock{Sprint Implementasi Workshop}{Project Based Learning (PBL) -- Sprint 1-2, UTS Milestone, Sprint 3-4}{SCPMK0802-05402, SCPMK0603-05401, SCPMK0803-05403}{Mahasiswa mengeksekusi empat sprint pengembangan produk teknologi terapan, mencakup implementasi fitur, pengujian, review stakeholder, dan milestone UTS, secara berkelanjutan hingga produk siap industri.}{Melaksanakan sprint planning, mengimplementasikan fitur sesuai backlog, melakukan pengujian (unit, integration, regression), mereview hasil bersama stakeholder, dan mendokumentasikan setiap sprint.}{Kode sumber teruji, laporan sprint, bukti demo kepada stakeholder, dokumentasi pengujian per sprint, dan evidence UTS milestone.}{Sprint planning dan backlog execution & 18 & Backlog sprint terencana dan dieksekusi dengan baik; velocity tim konsisten; task diselesaikan sesuai definition of done. & Backlog ada dan sebagian besar task diselesaikan, tetapi konsistensi velocity atau definition of done masih belum optimal. & Backlog tidak terkelola, task sering tidak selesai, atau tidak ada definisi done yang jelas. \\ \\
Kualitas implementasi produk per sprint & 20 & Fitur diimplementasikan sesuai kebutuhan industri, kode bersih dan terstruktur, serta setiap sprint menghasilkan increment yang dapat di-demo. & Implementasi berjalan tetapi kualitas kode atau kesesuaian dengan kebutuhan masih perlu perbaikan pada beberapa fitur. & Implementasi tidak sesuai kebutuhan, kode sulit dipelihara, atau sprint tidak menghasilkan increment yang dapat di-demo. \\ \\
Hasil pengujian dan integrasi testing & 17 & Test plan komprehensif dieksekusi, defect ditracking dan ditangani, dan integrasi antarmodul tervalidasi di setiap sprint. & Pengujian dilakukan tetapi cakupan atau tracking defect belum konsisten di semua sprint. & Pengujian tidak terstruktur, banyak defect tidak tertangani, atau integrasi antarmodul tidak divalidasi. \\}

\assessmentblock{Laporan Quality Assurance}{Case Method (CM) -- laporan teknis QA}{SCPMK0603-05401}{Mahasiswa menyusun laporan QA komprehensif yang mendokumentasikan seluruh aktivitas pengujian selama workshop, termasuk test plan, hasil pengujian, analisis defect, dan rekomendasi perbaikan produk.}{Mengumpulkan dan menganalisis data pengujian dari semua sprint, menyusun laporan defect, mengidentifikasi pola masalah kualitas, dan merumuskan rekomendasi perbaikan.}{Laporan QA lengkap, test summary report, defect log, matriks pengujian, dan rekomendasi peningkatan kualitas.}{Kelengkapan dokumentasi test plan dan hasil pengujian & 4 & Test plan mencakup semua jenis pengujian (unit, integration, UAT), hasil terdokumentasi lengkap, dan coverage terukur. & Test plan dan hasil tersedia tetapi cakupan pengujian atau kelengkapan dokumentasi masih terbatas. & Dokumentasi pengujian tidak lengkap, test plan tidak ada, atau coverage tidak dapat diukur. \\ \\
Analisis temuan dan defect tracking & 3 & Defect dianalisis berdasarkan severity/priority, pola masalah teridentifikasi, dan tindak lanjut setiap defect terdokumentasi. & Analisis defect cukup, tetapi klasifikasi atau tracking tindak lanjut masih belum komprehensif. & Defect tidak dianalisis secara sistematis atau tidak ada tracking tindak lanjut yang jelas. \\ \\
Rekomendasi perbaikan dan peningkatan kualitas & 3 & Rekomendasi spesifik, berbasis data temuan, dapat diterapkan, dan diprioritaskan berdasarkan dampak terhadap produk. & Rekomendasi relevan tetapi masih umum atau tidak sepenuhnya berbasis data temuan pengujian. & Rekomendasi tidak operasional, tidak terkait temuan, atau tidak memberikan nilai tambah bagi kualitas produk. \\}

\assessmentblock{UAS Demo Produk dan Defense}{UAS demo dan defense kepada industri}{SCPMK0803-05403, SCPMK0802-05402}{Mahasiswa mendemonstrasikan produk teknologi terapan hasil workshop kepada stakeholder industri dan melakukan defense atas keputusan teknis dan manajemen proyek yang diambil selama workshop berlangsung.}{Mempersiapkan demo produk, menyusun bahan presentasi, melakukan demo kepada panel industri, dan menjawab pertanyaan terkait keputusan teknis dan manajemen proyek.}{Demo produk yang berjalan, slide presentasi, laporan akhir workshop, dan dokumentasi proyek lengkap.}{Kualitas produk yang didemonstrasikan & 8 & Produk berjalan stabil, semua fitur utama terdemo dengan baik, dan kualitas memenuhi standar industri yang ditetapkan. & Produk berjalan untuk sebagian besar fitur, tetapi beberapa fitur masih memiliki kekurangan atau belum stabil. & Produk tidak berjalan dengan baik atau banyak fitur utama yang tidak dapat didemonstrasikan. \\ \\
Kemampuan defense dan penjelasan keputusan teknis & 7 & Mahasiswa menjelaskan keputusan arsitektur, teknologi, dan pengujian dengan argumentasi berbasis data dan pengalaman industri. & Defense cukup jelas, tetapi beberapa keputusan teknis belum didukung argumentasi yang kuat. & Tidak mampu menjelaskan keputusan teknis atau defense tidak mencerminkan pemahaman mendalam terhadap produk. \\ \\
Manajemen proyek dan kinerja tim workshop & 5 & Manajemen backlog, distribusi kerja tim, dan capaian sprint terdokumentasi dan dipertanggungjawabkan secara profesional. & Manajemen proyek cukup baik tetapi beberapa aspek koordinasi atau dokumentasi kinerja tim masih kurang lengkap. & Manajemen proyek tidak terstruktur atau tim tidak dapat mempertanggungjawabkan proses workshop secara kolektif. \\}

\begin{tblr}{width=\textwidth,colspec={|Q[l,m,3.2cm]|X[l,m]|},hlines,vlines,hspan=minimal,row{1-3}={bg=headgray},rowsep=1.5pt,colsep=3pt}
Jadwal Pelaksanaan: & Minggu 1--17 sesuai RPS. Setiap evaluasi memiliki RTM dan rubrik multi-kriteria. \\
Lain-Lain Yang Diperlukan: & LMS, repositori kode sumber, tools CI/CD dan testing, perangkat lunak sesuai mata kuliah, dan perangkat presentasi. \\
Pustaka: & Mengacu pustaka utama dan pendukung pada bagian RPS. \\
\end{tblr}
